\section{Plano de formalização e conclusão}
\label{sec:plano}
A ordem abaixo organiza as dependências de abertura e lançamento. Não representa cronograma executado, atribuição individual já aceita ou prazo oficial de órgãos públicos.

{\footnotesize
\renewcommand{\arraystretch}{1.0}
\begin{longtable}{@{}L{0.8cm} L{6.2cm} L{4.5cm} L{3.3cm}@{}}
\caption{Checklist proposto para a formalização do PetCuida.}\\
\toprule
\cabfundo \cab{Nº} & \cab{Ação e evidência esperada} & \cab{Responsável proposto} & \cab{Momento}\\
\midrule
\endfirsthead
\toprule
\cabfundo \cab{Nº} & \cab{Ação e evidência esperada} & \cab{Responsável proposto} & \cab{Momento}\\
\midrule
\endhead
\bottomrule
\endfoot
1 & Confirmar fundadores, quotas, capital e escopo real do produto. Evidência: decisões documentadas. & Grupo fundador. & Antes do contrato.\\
2 & Pesquisar nome e marca; realizar viabilidade de endereço e atividades. Evidência: consultas e análise. & Grupo e contador; apoio jurídico para marca. & Antes do registro.\\
3 & Redigir contrato social, acordo de sócios e cessões de direitos. Evidência: instrumentos aprovados. & Fundadores e assessoria jurídica. & Antes da abertura.\\
4 & Registrar na JUCEMG e obter CNPJ pelo fluxo integrado; verificar inscrições e licenças. & Contador e administrador indicado. & Abertura.\\
5 & Comparar regimes fiscais e formalizar a opção cabível no prazo vigente. Evidência: simulação e protocolo. & Contador, com aprovação dos sócios. & Na abertura.\\
6 & Abrir conta empresarial, organizar escrituração, emissão fiscal e conciliação de repasses. & Administração e contador. & Antes de movimentar recursos.\\
7 & Confirmar exigências do CRMV, limites da orientação e documentação das clínicas. & Assessoria jurídica e médico-veterinário habilitado. & Antes de ofertar serviços.\\
8 & Implantar termos, privacidade, segurança, contratos de parceiros e regras dos créditos. & Equipe de produto e assessoria especializada. & Antes do lançamento.\\
9 & Depositar pedido de marca e avaliar registro de software e licenças. & Administração e titulares dos direitos. & Antes de ampliar divulgação.\\
10 & Revisar contratos, tributação, licenças e governança quando houver novas receitas ou investimento. & Administração, contador e assessoria jurídica. & Durante a operação.\\
\end{longtable}}

\subsection{Síntese da decisão}
A \textbf{LTDA com quatro sócios} apresenta a melhor adequação ao \projeto{} nesta fase: admite a equipe fundadora, permite divisão de quotas e oferece estrutura de responsabilidade e governança mais compatível com o risco do negócio do que MEI ou EI, sem a complexidade inicial de uma S.A. A SLU permanece possível apenas em cenário de fundador único.

A formalização combina registro, CNPJ, licenciamento, tributação, contratos, proteção intelectual e LGPD. Para o PetCuida, destacam-se os limites da atuação veterinária e os repasses de créditos. Sócios e quotas propostos devem ser confirmados; o Simples depende de análise contábil. As decisões exigem revisão profissional antes da abertura.